\documentclass[12pt,a4paper]{article}
\usepackage[utf8]{inputenc}
\usepackage[T1]{fontenc}
\usepackage[ngerman]{babel}
\usepackage{geometry}
\usepackage{hyperref}
\usepackage{listings}
\usepackage{xcolor}
\usepackage{booktabs}
\usepackage{graphicx}
\usepackage{amsmath}
\usepackage{enumitem}
\usepackage{titlesec}
\usepackage{fancyhdr}

\geometry{left=2.5cm, right=2.5cm, top=3cm, bottom=3cm}

\pagestyle{fancy}
\fancyhf{}
\rhead{MATSE-Eisenbahngesellschaft -- Fahrplansystem}
\lhead{IHK Abschlussprüfung Sommer 2026}
\rfoot{Seite \thepage}

\lstset{
  basicstyle=\ttfamily\small,
  keywordstyle=\color{blue},
  commentstyle=\color{gray},
  stringstyle=\color{red},
  breaklines=true,
  frame=single,
  language=Java,
  showstringspaces=false,
}

% -----------------------------------------------------------------------
\begin{document}

\begin{titlepage}
  \centering
  \vspace*{2cm}
  {\LARGE\bfseries Entwicklung eines Mathematisch-technischen Softwaresystems\\[0.5em]}
  {\large IHK Abschlussprüfung Sommer 2026 -- 6511\\[1em]}
  {\large\bfseries Fahrplanoptimierung für eingleisige Eisenbahnstrecken\\[2em]}
  {\large MATSE-Eisenbahngesellschaft\\[3em]}
  \vfill
  {\normalsize Abgabedatum: \today}
\end{titlepage}

\tableofcontents
\newpage

% -----------------------------------------------------------------------
\section{Aufgabenanalyse und Verfahrensbeschreibung}

\subsection{Aufgabenanalyse}

Die MATSE-Eisenbahngesellschaft betreibt eingleisige Strecken, auf denen sich
entgegenkommende Züge nur an Bahnhöfen begegnen dürfen. Ziel ist die Berechnung
eines getakteten Fahrplans, der:
\begin{itemize}
  \item Kollisionen zwischen Hin- und Rückfahrt-Zügen verhindert,
  \item Wartezeiten minimiert und möglichst gleichmäßig auf beide Richtungen verteilt,
  \item einen 60-Minuten-Takt einhält (jeder Zug fährt stündlich am selben Bahnhof ab),
  \item eine Haltezeit von 1\,Minute je Zwischenbahnhof sowie eine Sicherheitswartezeit
        von 1\,Minute nach Einfahrt eines Gegenzugs berücksichtigt.
\end{itemize}

\noindent
Es sind drei Strategien zu implementieren:
\begin{enumerate}
  \item \textbf{Einfache Fahrt} -- Referenzplan ohne Kollisionsvermeidung.
  \item \textbf{Einseitiges Warten} -- Rückfahrt wartet auf freien Streckenabschnitt.
  \item \textbf{Beidseitiges Warten} -- Wartezeiten werden auf beide Richtungen aufgeteilt.
\end{enumerate}

% -----------------------------------------------------------------------
\subsection{Entwurf des Einlesens und der Initialisierung der Daten}

Die Eingabedatei besteht aus drei Abschnitten:
\begin{itemize}
  \item \texttt{Strecke:} -- durch Whitespace getrennte Bahnhofsnamen (eine Zeile).
  \item \texttt{Abstaende:} -- Fahrzeiten (Minuten) zwischen aufeinanderfolgenden
        Bahnhöfen (eine Zeile mit $n-1$ Werten für $n$ Bahnhöfe).
  \item \texttt{Start Hinfahrt:} -- Abfahrtszeit der Hinfahrt am ersten Bahnhof (reine Minutenangabe).
\end{itemize}

\noindent
Die Klasse \texttt{EingabeLeser} liest die Datei zeilenweise ein, erkennt die
Schlüsselwörter und befüllt ein \texttt{BahnhofPlan}-Objekt mit den Bahnhöfen
und Abständen sowie die Startzeit für die Simulation.

% -----------------------------------------------------------------------
\subsection{Kollisionserkennung und Wartezeitermittlung}

Eine Kollision liegt vor, wenn sich Hin- und Rückfahrtzug gleichzeitig auf
demselben Streckenabschnitt (zwischen zwei Bahnhöfen) befinden. \\[0.3em]

\noindent
Sei $[h_{\text{start}},\, h_{\text{ende}}]$ das Zeitfenster des Hinfahrtzugs und
$[r_{\text{start}},\, r_{\text{ende}}]$ das des Rückfahrtzugs auf einem Abschnitt.
Eine Kollision tritt auf, wenn sich diese Intervalle im 60-Minuten-Takt
\emph{überschneiden}. Da Zeiten modulo\,60 gerechnet werden, prüft die Methode
\texttt{isKollision()} für jede Minute $m \in [0,59]$, ob $m$ in beiden
Intervallen liegt.

\bigskip
\noindent
Die zusätzliche Wartezeit an einem Bahnhof $B$ berechnet sich als:
\[
  w_B = \bigl(\text{Ankunft}_{\text{Hin}}(B) - \text{Ankunft}_{\text{Rück}}(B) + 60\bigr) \bmod 60
\]

% -----------------------------------------------------------------------
\subsection{Entwurf eines Algorithmus zur Durchführung der Strategie „Einfache Fahrt"}

\begin{itemize}
  \item Die Hinfahrt startet zur gegebenen Startzeit und fährt an jedem Bahnhof
        nach genau 1\,Minute Halt schnellstmöglich weiter.
  \item Die Rückfahrt startet unmittelbar nach Ankunft am Endbahnhof
        (Ankunftszeit + 1\,Minute Halt).
  \item Es wird \emph{keinerlei} Rücksicht auf den Folgeabschnitt genommen.
  \item Kollisionen werden erkannt und in der Ausgabe mit \texttt{x} markiert;
        der Plan wird trotzdem ausgegeben (als Referenz).
  \item Wartezeiten: $w_{\text{Hin}} = 0$, $w_{\text{Rück}} = 0$,
        Score $= 0^2 + 0^2 = 0$.
\end{itemize}

% -----------------------------------------------------------------------
\subsection{Entwurf eines Algorithmus zur Durchführung der Strategie „Einseitiges Warten"}

\begin{itemize}
  \item Die \textbf{Hinfahrt} wird identisch zu „Einfache Fahrt" berechnet
        (schnellstmögliche kollisionsfreie Verbindung).
  \item Die \textbf{Rückfahrt} startet direkt nach Ankunft des Hinfahrtzugs am
        Endbahnhof (inkl. Sicherheitswartezeit von 1\,Minute).
  \item Vor jedem Abfahren vom aktuellen Bahnhof $B$ in Richtung $B-1$ wird
        geprüft, ob der Streckenabschnitt $[B, B-1]$ frei ist:
        \begin{enumerate}
          \item Falls \emph{keine} Kollision droht: sofortige Abfahrt.
          \item Falls eine Kollision droht: Warten, bis der Gegenzug in $B$
                eingetroffen ist, dann erst Abfahrt (nach 1\,Minute Sicherheitswartezeit).
        \end{enumerate}
  \item Die Wartezeiten kumulieren sich ausschließlich auf der Rückfahrt.
  \item Score $= 0^2 + \left(\sum w_{\text{Rück}}\right)^2$.
\end{itemize}

% -----------------------------------------------------------------------
\subsection{Entwurf eines Algorithmus zur Durchführung der Strategie „Beidseitiges Warten"}

\begin{itemize}
  \item Ausgangspunkt ist das Ergebnis der Strategie „Einseitiges Warten"
        mit den berechneten Wartezeiten je Bahnhof auf der Rückfahrt.
  \item Für jeden Bahnhof $B$ mit $w_{\text{Rück}}(B) > 1$ wird die Wartezeit
        gleichmäßig aufgeteilt:
        \[
          w_{\text{Hin}}(B) = \left\lfloor \frac{w_{\text{Rück}}(B)}{2} \right\rfloor,
          \quad
          w_{\text{Rück}}(B) \mathrel{-}= w_{\text{Hin}}(B)
        \]
  \item Die Abfahrtszeit der Rückfahrt am Endbahnhof wird um
        $\Delta = \sum_B w_{\text{Hin}}(B)$ nach vorne verschoben.
  \item Anschließend werden Hin- und Rückfahrt mit den aktualisierten
        Wartezeiten neu berechnet (\texttt{hinfahrtMitKollision},
        \texttt{rueckfahrtMitKollision}).
  \item Score $= \left(\sum w_{\text{Hin}}\right)^2 + \left(\sum w_{\text{Rück}}\right)^2$.
\end{itemize}

% -----------------------------------------------------------------------
\subsection{Entwurf der Ausgabe gemäß Aufgabenstellung}

Die Ausgabe erfolgt je Strategie in tabellarischer Form:
\begin{enumerate}
  \item Ankunftszeiten der Hinfahrt (\texttt{An})
  \item Zusätzliche Wartezeiten der Hinfahrt (\texttt{Wa}) -- leer falls 0
  \item Abfahrtszeiten der Hinfahrt (\texttt{Ab})
  \item Bahnhofsnamen -- mit \texttt{x} zwischen zwei Bahnhöfen bei Kollision
  \item Abfahrtszeiten der Rückfahrt (\texttt{Ab})
  \item Zusätzliche Wartezeiten der Rückfahrt (\texttt{Wa}) -- leer falls 0
  \item Ankunftszeiten der Rückfahrt (\texttt{An})
  \item Kennwerte: Gesamtdauer, Summe Wartezeiten, Score (Summe Strafen)
\end{enumerate}

% -----------------------------------------------------------------------
\section{Objektorientierte Programmkonzeption}

\subsection{Entwurf der Klassen, Methoden und Datenstrukturen}

Das Programm nutzt eine doppelt verkettete Liste aus \texttt{BahnhofNode}-Objekten,
die im \texttt{BahnhofPlan} verwaltet wird. Folgende Klassen wurden entworfen:

\begin{itemize}
  \item \textbf{BahnhofNode} -- Knoten der Liste; speichert Name, Ankunfts-/Abfahrtszeiten
        für Hin- und Rückfahrt, Wartezeiten und Abstand zum nächsten Bahnhof.
  \item \textbf{BahnhofPlan} -- doppelt verkettete Liste der Bahnhöfe; stellt
        \texttt{head}/\texttt{tail} sowie Methoden zum Aufbau und Reset bereit.
  \item \textbf{FahrplanStrategie} (abstrakt) -- definiert die Schnittstelle aller
        Strategien; enthält die gemeinsamen statischen Berechnungsmethoden
        (\texttt{hinfahrtOhneKollision}, \texttt{rueckfahrtMitKollision}, etc.).
  \item \textbf{EinfacheFahrtStrategie} -- implementiert „Einfache Fahrt".
  \item \textbf{EinseitigesWartenStrategie} -- implementiert „Einseitiges Warten".
  \item \textbf{BeidseitigesWartenStrategie} -- implementiert „Beidseitiges Warten".
  \item \textbf{EingabeLeser} -- liest und parst die Eingabedatei.
  \item \textbf{AusgabeManager} -- formatiert und druckt die Ergebnistabellen.
  \item \textbf{SimulationsDaten} (Record) -- Datencontainer für Plan und Startzeit.
  \item \textbf{Main} -- Einstiegspunkt; orchestriert Einlesen, Strategieauswahl und Ausgabe.
\end{itemize}

% -----------------------------------------------------------------------
\subsection{UML-Sequenzdiagramme für die wesentlichen Abläufe}

Die folgenden wesentlichen Abläufe sind als UML-Sequenzdiagramme zu dokumentieren
(Diagramme in separater Anlage):

\begin{itemize}
  \item Programmstart: \texttt{Main} $\to$ \texttt{EingabeLeser} $\to$ \texttt{BahnhofPlan}
  \item Strategieausführung: \texttt{Main} $\to$ \texttt{FahrplanStrategie}
        $\to$ \texttt{BahnhofNode} (Kollisionsprüfung)
  \item Ausgabe: \texttt{Main} $\to$ \texttt{AusgabeManager} $\to$ \texttt{BahnhofNode}
\end{itemize}

% -----------------------------------------------------------------------
\subsection{Detaillierte Beschreibung der wesentlichen Methoden}

Die folgenden Methoden sind in Form von Nassi-Shneiderman-Diagrammen oder
Programmablaufplänen detailliert beschrieben (Diagramme in separater Anlage):

\begin{itemize}
  \item \texttt{FahrplanStrategie.rueckfahrtMitKollision()} -- Kernlogik von
        „Einseitiges Warten": iteriert rückwärts über den Plan, prüft Kollision,
        setzt Wartezeiten.
  \item \texttt{BahnhofNode.isKollision()} -- prüft Intervallüberschneidung im
        60-Minuten-Raster für einen Streckenabschnitt.
  \item \texttt{FahrplanStrategie.beidseitigesWarten()} -- verteilt Wartezeiten
        hälftig auf beide Richtungen und berechnet den Plan neu.
  \item \texttt{AusgabeManager.druckeErgebnis()} -- erzeugt die tabellarische
        Ausgabe je Strategie.
\end{itemize}

% -----------------------------------------------------------------------
\section{Dokumentation}

\subsection{Verbale Beschreibung der Kollisionserkennung, Wartezeitermittlung
            und der drei Strategien}

\subsubsection{Kollisionserkennung}
Zwei Züge kollidieren auf einem Streckenabschnitt $[B_i, B_{i+1}]$, wenn sich
ihre Zeitfenster im Takt von 60\,Minuten überlappen. Die Methode
\texttt{isKollision()} in \texttt{BahnhofNode} prüft für jeden Wert
$m \in \{0,\ldots,59\}$, ob $m$ sowohl im Hinfahrt-Intervall
$[h_{\text{dep}}(B_i),\, h_{\text{arr}}(B_{i+1})]$ als auch im
Rückfahrt-Intervall $[r_{\text{dep}}(B_{i+1}),\, r_{\text{arr}}(B_i)]$ liegt.
Dabei werden Intervalle, die den Stundenwechsel überschreiten (z.\,B.
$[55, 05]$), korrekt behandelt.

\subsubsection{Wartezeitermittlung}
Muss ein Rückfahrtzug in Bahnhof $B$ warten, so berechnet sich die Wartezeit als
Differenz zwischen der Ankunft des Gegenzugs und der eigenen (frühestmöglichen)
Abfahrt, zuzüglich der Sicherheitswartezeit von 1\,Minute.

\subsubsection{Einfache Fahrt -- am Beispiel}
Eingabe (7 Bahnhöfe, Abstände 5\,8\,6\,7\,3\,2, Start 17\,min):
Die Hinfahrt fährt: $17 \xrightarrow{5} 22 \xrightarrow{+1} 23 \xrightarrow{8}
31 \xrightarrow{+1} \ldots$
Die Rückfahrt startet sofort nach Ankunft. Zwischen B und C tritt eine Kollision
auf ($\to$ Markierung mit \texttt{x} in der Ausgabe).

\subsubsection{Einseitiges Warten -- am Beispiel}
Die Hinfahrt ist identisch. Am Bahnhof C (Rückfahrt) wird eine drohende Kollision
erkannt: der Rückfahrtzug wartet 16\,Minuten. Score $= 16^2 = 256$.

\subsubsection{Beidseitiges Warten -- am Beispiel}
Die 16\,Warteminuten aus der Rückfahrt werden nicht verwendet;
durch eine andere Wahl der Rückfahrt-Startzeit lässt sich die Begegnung
in Bahnhof E verlegen, ohne zusätzliche Wartezeit. Score $= 0$.

% -----------------------------------------------------------------------
\subsection{Entwicklerdokumentation und Benutzeranleitung}

\subsubsection{Entwicklerdokumentation}

\paragraph{Voraussetzungen}
\begin{itemize}
  \item Java 17 oder höher
  \item Apache Maven 3.8 oder höher
\end{itemize}

\paragraph{Build}
\begin{lstlisting}[language=bash]
mvn compile
\end{lstlisting}

\paragraph{Ausführung}
\begin{lstlisting}[language=bash]
# Aus dem übergeordneten Verzeichnis aufrufen,
# damit der relative Pfad zur Datei stimmt:
mvn exec:java -Dexec.mainClass="de.matse.ihk.Main"
\end{lstlisting}

\paragraph{Projektstruktur}
\begin{lstlisting}
eisenbahn-fahrplan/
  src/main/java/de/matse/ihk/
    BahnhofNode.java           -- Listenknoten
    BahnhofPlan.java           -- Doppelt verkettete Liste
    FahrplanStrategie.java     -- Abstrakte Basisklasse
    EinfacheFahrtStrategie.java
    EinseitigesWartenStrategie.java
    BeidseitigesWartenStrategie.java
    SimulationsDaten.java      -- Record (Eingabe-Container)
    Main.java                  -- Einstiegspunkt
    io/
      EingabeLeser.java
      AusgabeManager.java
  datei.txt                    -- Aktuelle Eingabedatei
\end{lstlisting}

\subsubsection{Benutzeranleitung}

\begin{enumerate}
  \item Eingabedatei \texttt{datei.txt} im Projektverzeichnis anlegen/bearbeiten
        (Format: siehe Abschnitt 1.2).
  \item Programm starten (siehe oben).
  \item Die Ausgabe erscheint in der Konsole; sie enthält Fahrpläne und Kennwerte
        für alle drei Strategien.
\end{enumerate}

\paragraph{Eingabeformat}
\begin{lstlisting}
Strecke:
A B C ...

Abstaende:
<d1> <d2> ...

Start Hinfahrt:
<startminute>
\end{lstlisting}

% -----------------------------------------------------------------------
\subsection{Zusammenfassung mit Ausblick}

Das implementierte System berechnet getaktete Fahrpläne für eingleisige
Eisenbahnstrecken und vergleicht drei Strategien anhand eines quadratischen
Warte\-zeit-Scores.

\medskip
\noindent
\textbf{Erreichte Ziele:}
\begin{itemize}
  \item Alle drei Strategien korrekt implementiert und getestet.
  \item Kollisionserkennung mit Behandlung des Stundenwechsels.
  \item Tabellarische Ausgabe gemäß Aufgabenstellung.
  \item Score-Berechnung nach quadratischem Strafmaß.
\end{itemize}

\noindent
\textbf{Ausblick und mögliche Erweiterungen:}
\begin{itemize}
  \item Unterstützung variabler Haltezeiten und Sicherheitswartezeiten
        (bereits als Konstante vorbereitet, leicht anpassbar).
  \item Grafische Ausgabe des Fahrplans (z.\,B. als Bildfahrplan).
  \item Automatische Optimierung der Rückfahrt-Startzeit für
        „Beidseitiges Warten" durch Brute-Force über alle 60\,möglichen
        Startzeiten.
  \item Einlesen mehrerer Strecken aus einer Datei.
\end{itemize}

% -----------------------------------------------------------------------
\section{Diskussion der Testbeispiele}

\subsection{Skript zur automatischen Ausführung der Beispiele}

\begin{lstlisting}[language=bash]
#!/bin/bash
# run_tests.sh -- alle Testfaelle automatisch ausfuehren
EXAMPLES_DIR="docs/testbeispiele"
for f in "$EXAMPLES_DIR"/*.txt; do
  echo "=== $f ==="
  cp "$f" datei.txt
  mvn -q exec:java -Dexec.mainClass="de.matse.ihk.Main"
done
\end{lstlisting}

% -----------------------------------------------------------------------
\subsection{Angegebene Beispiele}

\subsubsection{Beispiel 1 -- 3 Bahnhöfe (einfach, keine Wartezeit)}

\textbf{Eingabe:}
\begin{lstlisting}
Strecke:
A B C
Abstaende:
5 7
Start Hinfahrt:
17
\end{lstlisting}

\textbf{Erwartetes Ergebnis:} Alle drei Strategien identisch, Score $= 0$.
Keine Kollision, da die Strecke kurz genug ist.

\subsubsection{Beispiel 2 -- 7 Bahnhöfe (ungerade, Wartezeit bei Einseitigem Warten)}

\textbf{Eingabe:}
\begin{lstlisting}
Strecke:
A B C D E F G
Abstaende:
5 8 6 7 3 2
Start Hinfahrt:
17
\end{lstlisting}

\textbf{Einfache Fahrt:} Kollision zwischen B und C (Markierung \texttt{x}).

\textbf{Einseitiges Warten:}
\begin{lstlisting}
An    22 31 38 46 50 53
Wa
Ab 17 23 32 39 47 51 54
    A  B  C  D  E  F  G
Ab 31 41 32 09 01 57 54
Wa       16
An 46 40 15 08 00 56
Gesamtdauer Hinfahrt, Rueckfahrt       : 36, 52
Summe Wartezeiten Hinfahrt, Rueckfahrt : 0, 16
Summe Strafen                          : 256
\end{lstlisting}

Der Rückfahrtzug muss in C \textbf{16 Minuten} warten (ungerade Anzahl von
Bahnhöfen führt zu einer asymmetrischen Streckenaufteilung).

\textbf{Beidseitiges Warten:} Score $= 0$ durch optimale Wahl der Rückfahrt-Startzeit.

\subsubsection{Beispiel 3 -- 10 Bahnhöfe (gerade, mehrfache Begegnungen)}

\textbf{Eingabe:}
\begin{lstlisting}
Strecke:
A B C D E F G H I J
Abstaende:
12 8 14 3 5 21 13 6 8
Start Hinfahrt:
17
\end{lstlisting}

\textbf{Einseitiges Warten:} Gesamtwartezeit Rückfahrt $= 8\,\text{min}$,
Score $= 64$.\\
\textbf{Beidseitiges Warten:} Wartezeit je Richtung $= 4\,\text{min}$,
Score $= 4^2 + 4^2 = 32$.

% -----------------------------------------------------------------------
\subsection{Weitere repräsentative Beispiele mit Begründung der Wahl}

\subsubsection{Beispiel 4 -- 2 Bahnhöfe (minimale Strecke)}

\textbf{Begründung:} Grenzfall mit nur einem Streckenabschnitt; überprüft, ob
die Schleifenlogik korrekt abbricht und keine Division durch null auftritt.

\begin{lstlisting}
Strecke:
A B
Abstaende:
15
Start Hinfahrt:
0
\end{lstlisting}

Erwartetes Ergebnis: Keine Wartezeit, Score $= 0$, da kein Zwischenbahnhof
existiert, an dem Züge aufeinandertreffen könnten.

\subsubsection{Beispiel 5 -- Große Fahrtzeit (Kollision erzwungen)}

\textbf{Begründung:} Strecke mit sehr langen Teilabschnitten; testet, ob
Wartezeiten korrekt über den Stundenwechsel (Modulo\,60) berechnet werden.

\begin{lstlisting}
Strecke:
A B C
Abstaende:
29 29
Start Hinfahrt:
0
\end{lstlisting}

% -----------------------------------------------------------------------
\subsection{Grenz- und Fehlerfälle}

\begin{itemize}
  \item \textbf{Nur 2 Bahnhöfe:} Keine Haltezeiten an Zwischenbahnhöfen;
        Mindestdauer $=$ Abstand.
  \item \textbf{Startzeit 0:} Erlaubt; Zeiten bleiben nicht-negativ durch
        Modulo\,60-Rechnung.
  \item \textbf{Abstand = 30\,min (halber Takt):} Hin- und Rückfahrtzug erreichen
        den Mittelpunkt gleichzeitig -- maximale Wartezeit.
  \item \textbf{Ungültige Eingabedatei:} \texttt{EingabeLeser} wirft
        \texttt{IOException}; \texttt{Main} gibt eine Fehlermeldung aus
        und beendet das Programm.
  \item \textbf{Einzelner Bahnhof:} Kein sinnvoller Fahrplan möglich; wird
        nicht explizit abgefangen (Vorbedingung: mindestens 2 Bahnhöfe).
\end{itemize}

\end{document}
